\section{Supplementary Learner Selection}
\label{sec:selection}
Forecasting a censored requirement and choosing a learner at one budget have different losses. We therefore treat selection as a secondary experiment. The final rule fits two ExtraTrees regressors on the opened fourth cohort's selection outcomes to predict each of the seven registered learners' seed success rate and normalized return at $x10$. Its inputs combine CFQI pilot features, CFQI's pre-outcome \OEH{} forecast probabilities, and a one-hot synthetic MDP-family label available in this study. Each regressor has 150 trees, minimum leaf size seven, and maximum feature fraction $0.8$. The rule maximizes predicted success plus $0.15$ times predicted return, but switches away from fixed CFQI only when the composite-score advantage exceeds $0.1$. Choices for the new cohort were sealed before its learner outcomes. The family label is an additional study-specific input, so this selector is not presented as a pilot-only forecast comparison.

The selection comparators are all seven fixed learners, random choice, the earlier OEH-v5 router, and a group oracle used only as an upper bound. The v5 router addresses the same decision target and was registered, so its result is retained even though the paper's primary contribution is forecasting. Fixed learner implementations include project tabular CFQI, IQL, and behavior cloning, and recorded author-code or library implementations of BCQ, CQL, SPIBB, and VI-LCB \cite{fujimoto2019bcq,kumar2020cql,kostrikov2022iql,laroche2019spibb,li2024sample}. This is a comparison of the registered implementations at matched group identities, not a claim about every possible tuning of the original algorithms.
